\chapter{Justificativa}
\pagestyle{plain}

A adoção acelerada da Inteligência Artificial por empresas públicas e privadas tem modificado profundamente os processos produtivos, a engenharia de software, a análise de dados, a automação industrial e a tomada de decisão estratégica. Tecnologias como IA Generativa, modelos de linguagem de larga escala (LLMs), agentes inteligentes, visão computacional e aprendizado de máquina passaram a integrar soluções utilizadas em setores como manufatura, saúde, agronegócio, finanças, logística, educação e serviços.

Entretanto, a velocidade dessa transformação tem evidenciado uma lacuna significativa na formação de profissionais capazes de compreender e integrar essas tecnologias em ambientes produtivos. Embora existam cursos voltados para temas específicos, como Ciência de Dados, Machine Learning ou Desenvolvimento de Software, ainda são escassas as formações que abordam de maneira integrada todo o ciclo de vida de uma solução de Inteligência Artificial.

Nesse contexto, este curso foi estruturado para desenvolver competências que abrangem desde os fundamentos da programação e da ciência de dados até a implantação e operação de sistemas inteligentes em produção, contemplando práticas modernas como MLOps, desenvolvimento assistido por IA, agentes inteligentes, RAG, computação distribuída e integração com dispositivos IoT.
